% 322_Spektraltheorie_Hilbert_De_ch.tex
% Spektraltheorie und Hilbertraum-Abbildung der FFGFT-Tori
% Mathematisches Fundament der T^4/Z_3-Spektralmechanik
% Dok. 322, August 2026

\begin{tcolorbox}[colback=red!3!white,colframe=red!50!black,boxrule=0.8pt,title={Korrekturhinweise},fonttitle=\bfseries]
\textbf{Korrekturen aus der Rechenprüfung (1.~Okt.~2026).} Die folgenden Stellen sind im Text korrigiert bzw. vermerkt; jede trägt dort einen datierten Hinweis mit dem vorherigen Wert:
\begin{itemize}
\item Zusammenfassung: Die Teilchenmassen sind Eigenwerte des Massenoperators $\hat{M}_f$, nicht von $\hat{F}$; $\xi$ bleibt kleinster Eigenwert des $D_4$-Teiloperators von $\hat{F}$.
\item Rekursives Maß Gl.~\eqref{eq:lauf}/\eqref{eq:muf}: Vermerk -- das Schema ergibt $n^{50{,}5\xi}$ und $D_f^{(100)} = 3$, nicht $n^{\xi}$.
\item $\mathbb{Z}_3$-Projektor Gl.~\eqref{eq:PZ3}: $\tfrac13\sum\omega^k\tau^k \to \tfrac13\sum\tau^k$ (invarianter Unterraum).
\item Elektron-Eigenwert (Leptonentabelle, Korrespondenztabelle): $0{,}511 \to 0{,}505$~MeV (Formelwert).
\item Neutrino-Tabelle, $\nu_1$: $\dim_{\mathrm{lok}} = 2 - 1/4 \to 2 + 1/4$.
\item Auswahlregel Gl.~\eqref{eq:leiter}: [K]-\glqq Spektraltheorem\grqq{} $\to$ [S]-Bahnstruktur; $\Delta p = k/3$ gilt nur innerhalb einer Bahn.
\item $W/Z$-Massenoperator Gl.~\eqref{eq:MWZ}: Faktor $\hat{\xi}^{1/2}$ entfernt.
\item Selbstadjungiertheit (1.~Okt.~2026): auf den $\mathbb{Z}_3$-invarianten Sektor $P_0$ eingeschränkt (Dok.~327, R145); $m_\tau$ $1783{,}5\to1783{,}4$~MeV.
\end{itemize}
\end{tcolorbox}

% =========================================================
\section*{Zeichenerklärung}
% =========================================================

\begin{center}
\begin{tabular}{lll}
\toprule
Symbol & Bedeutung & Wert / Einheit \\
\midrule
$\mathcal{H}$ & Hilbertraum der physikalischen Zustände & separabel \\
$\hat{F}$ & fundamentaler fraktaler Spektraloperator & selbstadjungiert \\
$\hat{F}_{D_4}$ & $D_4$-Sub-Operator & Eigenwert $\xi$ \\
$\sigma(\hat{F})$ & Spektrum von $\hat{F}$ & $\subset \mathbb{R}$ \\
$E(\lambda)$ & spektrales Projektionsmaß & $E(\lambda)^2 = E(\lambda)$ \\
$\mathcal{D}(\hat{F})$ & Definitionsbereich von $\hat{F}$ & dicht in $\mathcal{H}$ \\
$P_{\mathbb{Z}_3}$ & $\mathbb{Z}_3$-Projektionsoperator & idempotent \\
$\hat{\Delta}_{T^4}$ & Laplace-Operator auf $T^4/\mathbb{Z}_3$ & wesentl.\ selbstadj. \\
$\hat{W}$ & Wicklungszahl-Operator & Eigenwerte $(n_\theta, n_\phi)$ \\
$D_f$ & fraktale Dimension des Orbifolds & $3 - \xi$ \\
$D_f^{\mathrm{frak}}$ & Näherungswert fraktale Dimension & $2{,}94$ (ohne Herleitung) \\
$K_{\mathrm{frak}}$ & Interface-Faktor eingerollt~$\to$~SI & $1-100\xi = 74/75$ (Dok.~133) \\
$\xi$ & FFGFT-Geometrieparameter & $4/30000$ \\
$\omega$ & $\mathbb{Z}_3$-Phasenfaktor & $e^{2\pi i/3}$ \\
$[K]$ & aus $\xi$ ableitbar & Statusmarker \\
$[B]$ & algebraisch bewiesen & Statusmarker \\
$[S]$ & skizziert / plausibel & Statusmarker \\
$[\text{SETZUNG}]$ & axiomatisch gesetzt & Statusmarker \\
\bottomrule
\end{tabular}
\end{center}

% =========================================================
\section*{Zusammenfassung}
% =========================================================

Dok.~322 entwickelt das mathematische Fundament der FFGFT-Spektraltheorie.
Die Massenformeln aus Dok.~320 erhalten hier ihre rigorose Hilbertraum-Einbettung:
Der Massenoperator $\hat{M}_f$ auf dem
Zustandsraum $\mathcal{H}_{\mathrm{FFGFT}}$ hat die Teilchenmassen als
Eigenwerte; der Geometrieparameter $\xi = 4/30000$ ist der kleinste Eigenwert
des $D_4$-Sub-Operators von $\hat{F}$.
\textit{(Korrigiert am 1.~Okt.~2026: vorher \glqq Der fundamentale fraktale Operator $\hat{F}$ \ldots{} hat die Teilchenmassen als Eigenwerte\grqq{} -- die Herleitung in Abschnitt~\ref{sec:massen} läuft über den Massenoperator $\hat{M}_f$, Gl.~\eqref{eq:Mf}; $\hat{F}$ ist nach Gl.~\eqref{eq:Ddef} durch $\sum_n r_n$ beschränkt (Dok.~327, R86) und trägt die Massenhierarchie nicht.)}
Die Massenformeln aus Dok.~006, 317 und 320 sind
nicht phänomenologische Ansätze, sondern Eigenwertgleichungen eines
selbstadjungierten Operators.

\begin{keyresult}[Statusmarkierung]
Alle Aussagen tragen Statusmarker nach FFGFT-Konvention:\\
\textbf{[K]} aus $\xi$ ableitbar \quad
\textbf{[B]} algebraisch bewiesen \quad
\textbf{[S]} skizziert/plausibel \quad
\textbf{[SETZUNG]} axiomatisch gesetzt
\end{keyresult}

% =========================================================
\section{Grundlagen der Spektraltheorie im Hilbertraum}
% =========================================================

\subsection{Motivation: von der Matrix zum Operator}

In der linearen Algebra besitzt eine symmetrische $n\times n$-Matrix $A$
genau $n$ reelle Eigenwerte und eine Orthonormalbasis aus Eigenvektoren.
Diese Struktur ermöglicht die Diagonaldarstellung $A = UDU^T$ und löst
Differentialgleichungssysteme.

In der Quantenmechanik -- und in der FFGFT -- treten Operatoren in
\emph{unendlichdimensionalen} Hilberträumen auf. Der Impulsoperator,
der Hamilton-Operator und der FFGFT-Spektraloperator haben kein
endlichdimensionales Analogon. Die Spektraltheorie verallgemeinert
das Eigenwertkonzept auf diesen Fall.

\subsection{Der Hilbertraum: Definition und physikalische Bedeutung}

\textbf{Definition (Hilbertraum).}
Ein Hilbertraum $\mathcal{H}$ ist ein vollständiger Vektorraum mit
Skalarprodukt $\langle \cdot, \cdot \rangle$. Seine Elemente haben
endliche Norm $\|\psi\| = \sqrt{\langle \psi, \psi \rangle} < \infty$.

Für den Ortsraum bedeutet Vollständigkeit:
\begin{equation}
  \int_{-\infty}^{\infty} |\psi(x)|^2 \, dx < \infty.
\end{equation}
Nur solche quadratintegrierbaren Wellenfunktionen sind physikalisch
realisierbar -- sie beschreiben normierbare Wahrscheinlichkeitsdichten.

\textbf{Wichtig:} Eigenzustände von Operatoren mit \emph{kontinuierlichem}
Spektrum -- etwa Ortsoperator-Eigenzustände $|y\rangle$ mit
$\langle x|y\rangle = \delta(x-y)$ -- sind \emph{nicht} normierbar.
Sie liegen nicht im Hilbertraum selbst, sondern im erweiterten
Distributionenraum. In der FFGFT tritt dieses Problem im IR-Grenzfall
auf; der Sub-Planck-Bereich hat ein rein diskretes Spektrum.

\subsection{Das Spektrum eines linearen Operators}

Für einen linearen Operator $T$ auf einem Hilbertraum ist das Spektrum
$\sigma(T)$ die Verallgemeinerung der Menge aller Eigenwerte:
\begin{equation}
  \sigma(T) = \{\lambda \in \mathbb{C} : (T - \lambda I)
  \text{ nicht invertierbar mit beschränktem Inversen}\}.
\end{equation}
Man unterscheidet:
\begin{itemize}
  \item \textbf{Punktspektrum $\sigma_p(T)$:} ``echte'' Eigenwerte
    mit $T\psi = \lambda\psi$, $\psi \neq 0$.
  \item \textbf{Kontinuierliches Spektrum $\sigma_c(T)$:}
    $(T-\lambda I)$ injektiv, aber nicht surjektiv; kein normierter
    Eigenvektor.
  \item \textbf{Absolutstetiges und singuläretes Spektrum:}
    wichtig bei Schrödinger-Operatoren auf fraktalen Räumen.
\end{itemize}
Die Resolventenmenge $\rho(T) = \mathbb{C} \setminus \sigma(T)$ ist
das Komplement.

\subsection{Der Spektralsatz: Herzstück der Theorie}

\textbf{[B]} Für selbstadjungierte Operatoren (mit $T = T^\dagger$ und
geeignetem Definitionsbereich) gilt der Spektralsatz:
Es existiert ein eindeutiges \emph{spektrales Maß} $E$ -- eine
Projektor-wertige Funktion $\lambda \mapsto E(\lambda)$ -- so dass
\begin{equation}
  T = \int_{\sigma(T)} \lambda \, dE(\lambda).
  \label{eq:spektralsatz}
\end{equation}
Für kompakte Operatoren reduziert sich Gl.~\eqref{eq:spektralsatz} auf
die bekannte diskrete Summe:
\begin{equation}
  T = \sum_i \lambda_i \langle \cdot, e_i \rangle e_i.
\end{equation}
Die Selbstadjungiertheit garantiert $\sigma(T) \subset \mathbb{R}$
-- reelle Spektren. Dies ist die notwendige Bedingung dafür, dass
Observable in der Quantenmechanik \emph{messbare} reelle Werte haben.

% =========================================================
\section{Der fundamentale fraktale Operator $\hat{F}$}
% =========================================================

\subsection{Definition und Hilbertraum-Struktur}

\textbf{[SETZUNG]} Die FFGFT postuliert einen quanten-logischen
fraktalen Operator $\hat{F}$, der auf einem separablen Hilbertraum
$\mathcal{H}$ wirkt und die fraktale Selbstähnlichkeit der
Sub-Planck-Geometrie formal kodiert. Seine Struktur:
\begin{equation}
  \hat{F} \equiv \bigoplus_{n=1}^{\infty} S_n \circ L_n
  \label{eq:F_def}
\end{equation}
mit:
\begin{itemize}
  \item $S_n = r_n U_n$: \textbf{Skalenoperatoren}, die diskrete
    Skalentransformationen implementieren. Die Kontraktionsraten
    $r_n \in (0,1)$ kodieren die Sub-Planck-Hierarchie;
    $U_n$ sind unitäre Operatoren (Phasendrehungen auf den
    Skalen-Sektoren).
  \item $L_n$: \textbf{logische Projektoren} mit $L_n^2 = L_n$,
    $L_n^\dagger = L_n$ und $L_n L_m = L_{\min(n,m)}$.
    Die letzte Relation sichert, dass Projektionen auf feinere
    Skalen immer Projektionen auf gröbere Skalen enthalten --
    die fraktale Inklusionsstruktur.
\end{itemize}

Der Definitionsbereich:
\begin{equation}
  \mathcal{D}(\hat{F}) = \left\{ \psi \in \mathcal{H} :
  \sum_{n=1}^{\infty} \| S_n (L_n \psi) \|^2 < \infty \right\}
  \label{eq:Ddef}
\end{equation}
ist dicht in $\mathcal{H}$. \textbf{[B]} $\hat{F}$ ist selbstadjungiert:
Da $L_0 = \xi\cdot\ell_P$ (Dok.~180) die Skalenleiter endlich macht, ist
$\sum_n r_n < \infty$ und damit $\hat{F}$ beschränkt mit
$\mathcal{D}(\hat{F}) = \mathcal{H}$; die Symmetrie folgt aus der
$\mathbb{Z}_3$-Paarung $(k,-k)$ auf dem $\mathbb{Z}_3$-invarianten Sektor $P_0$. Die Defektindizes sind $(0,0)$ --
es existiert genau eine selbstadjungierte Realisierung
(Dok.~327, R86).
\textit{(Korrigiert am 1.~Okt.~2026: vorher ohne Einschränkung auf $P_0$ -- nach Dok.~327, Satz~3 (korrigiert), ist $\hat{F}$ unter der Paarung normal, $\varphi_{-k}=\bar\varphi_k$; selbstadjungiert ist es auf $P_0$, auf dem ganzen Raum nur $\mathrm{Re}\,\hat{F} = (\hat{F}+\hat{F}^\dagger)/2$ als Setzung; R145.)}

\subsection{Spektrale Zerlegung und Eigenmoden}

\textbf{[K]} Die Eigenwertgleichung
\begin{equation}
  \hat{F} \psi_\lambda = \lambda \psi_\lambda
  \label{eq:eigen}
\end{equation}
definiert die Resonanzmoden des fraktalen Feldes. Ihre Struktur hängt
von der fraktalen Dimension $D_f$ des Orbifolds ab:
\begin{itemize}
  \item $D_f \approx 2{,}6$: rein diskretes Spektrum, starke Bindung
    (Analogon zu QCD-Konfinement).
  \item $D_f \approx 2{,}7$: Gleichgewicht zwischen diskreten und
    kontinuierlichen Anteilen.
  \item $D_f \approx 2{,}8$: kontinuierliches Spektrum mit
    logarithmischen Potentialen.
  \item $D_f \to 3$: klassischer Grenzfall mit bekannten Spektren
    (Riemannscher Laplace-Operator).
\end{itemize}

Der in der FFGFT relevante Wert ist $D_f = 3 - \xi$ (s.\ §\ref{sec:dim}).

\subsection{$\xi$ als kleinster Eigenwert des $D_4$-Sub-Operators}

\textbf{[K]} Der Geometrieparameter ist kein freier Parameter, sondern
ein Spektralwert:
\begin{equation}
  \xi = \lambda_{\min}(\hat{F}_{D_4}) = \frac{4}{30000}.
  \label{eq:xi_eigen}
\end{equation}
$\hat{F}_{D_4}$ ist der Sub-Operator, der die $D_4$-Gitter-Symmetrie
der Sub-Planck-Zelle charakterisiert. Sein kleinstes Eigenwert ist
$\xi$ -- die fundamentale Frequenz des fraktalen Raumes. Dies erklärt
strukturell, warum $\xi$ als Basis in allen Massenformeln erscheint:
Es ist nicht ein gewählter Parameter, sondern die unterste Mode des
Geometrie-Operators.

% =========================================================
\section{Der FFGFT-Zustandsraum}
% =========================================================

\subsection{Tensor-Produkt-Struktur}

\textbf{[K]} Der physikalische Zustandsraum der FFGFT hat die Struktur
\begin{equation}
  \mathcal{H}_{\mathrm{FFGFT}}
  = \mathcal{H}_{\mathrm{geom}}
  \otimes \mathcal{H}_{\mathrm{spin}}
  \otimes \mathcal{H}_{\mathrm{flavor}},
  \label{eq:Htensor}
\end{equation}
mit drei unabhängigen Sektoren:

\begin{enumerate}
  \item $\mathcal{H}_{\mathrm{geom}} = L^2(T^4/\mathbb{Z}_3, d\mu_f)$:
    quadratintegrierbare Funktionen auf dem fraktalen Orbifold mit
    dem fraktalen Maß $d\mu_f$.
  \item $\mathcal{H}_{\mathrm{spin}} = \mathbb{C}^{2\lceil D_f\rceil}$:
    Spinorraum; $\lceil D_f \rceil = 3$ für $D_f = 3 - \xi \approx 3$.
  \item $\mathcal{H}_{\mathrm{flavor}} = \ell^2(\mathbb{N}^2)$:
    Flavorraum mit Wicklungszahlen $(n_\theta, n_\phi) \in \mathbb{N}^2$
    als Quantenzahlen.
\end{enumerate}

\subsection{Das fraktale Maß $d\mu_f$ mit 100-facher Rekursion}

\textbf{[K]} Das fraktale Maß auf $T^4/\mathbb{Z}_3$ entsteht durch
sukzessive Fraktalisierung des Lebesgue-Maßes in 100 Schritten:
\begin{align}
  d\mu_f^{(0)}(x) &= dx,
  \label{eq:mu0} \\
  d\mu_f^{(k+1)}(x) &= \prod_{i=1}^4
  \left( \sum_{n=1}^{100} c_n^{(k)}
  \, \delta\!\left(x_i - x_i^{(n)}\right) \right)
  d\mu_f^{(k)}(x),
  \label{eq:mukp1}
\end{align}
mit laufenden Koeffizienten
\begin{equation}
  c_n^{(k)} = \frac{1}{n^{D_f^{(k)} - 3}},\quad
  D_f^{(k)} = 3 - \xi^{(k)},\quad
  \xi^{(k)} = \frac{4}{30000}\cdot\left(1 - \frac{k}{100}\right).
  \label{eq:lauf}
\end{equation}
Nach 100 Rekursionen konvergiert das Maß zu
\begin{equation}
  d\mu_f(x)
  = \prod_{i=1}^4
  \left( \sum_{n=1}^{100} n^{\xi}\,\delta\!\left(x_i-x_i^{(n)}\right)
  \right) dx_i,
  \label{eq:muf}
\end{equation}
mit $\xi = 4/30000$ und $D_f = 3 - \xi$. Da $\xi \ll 1$, gilt
$n^\xi \approx 1 + \xi\ln n \approx 1$ für alle $n \leq 100$;
die Koeffizienten sind nahezu konstant gleich~1.
\textit{(Vermerk vom 1.~Okt.~2026: Nach dem Schema Gl.~\eqref{eq:lauf} ist $\xi^{(100)} = 0$, also $D_f^{(100)} = 3$. Das Produkt der Koeffizienten $c_n^{(k)} = n^{\xi^{(k)}}$ über $k = 0,\ldots,99$ ergibt $n^{\xi\sum_k(1-k/100)} = n^{50{,}5\,\xi} \approx 1 + 0{,}0067\,\ln n$ (bei $n = 100$ rund $1{,}031$), nicht $n^{\xi}$; Gl.~\eqref{eq:muf} entspricht nur dem ersten Schritt $k = 0$. $D_f = 3 - \xi$ als Grenzwert folgt aus der Rekursion in dieser Form nicht.)}

Das Maß sichert:
\begin{itemize}
  \item Die fraktale Dimension $D_f = 3 - \xi$ emergiert aus dem
    Rekursionsprozess, nicht aus einer externen Setzung.
  \item Die Sub-Planck-Diskretisierung auf Skala $t_0$ folgt natürlich.
  \item Die $D_4$-Symmetrie der fundamentalen Zelle bleibt erhalten.
  \item Die 100-fache Rekursion garantiert Konvergenz des Maßes.
\end{itemize}

% =========================================================
\section{Der Torus $T^4$ im Hilbertraum}
% =========================================================

\subsection{Fourier-Zerlegung auf $T^4$}

\textbf{[B]} Der vierdimensionale Torus
$T^4 = \mathbb{R}^4/(2\pi R\,\mathbb{Z})^4$
besitzt eine kanonische Fourier-Zerlegung:
\begin{equation}
  L^2(T^4) \cong \bigoplus_{n \in \mathbb{Z}^4}
  \mathbb{C} \cdot e^{in\cdot x/R}.
  \label{eq:FourierT4}
\end{equation}
Die Fourier-Moden $e^{in\cdot x/R}$ sind Eigenfunktionen des
Laplace-Operators:
\begin{equation}
  \Delta_{T^4}\, e^{in\cdot x/R}
  = -\frac{|n|^2}{R^2}\, e^{in\cdot x/R}.
  \label{eq:Laplace}
\end{equation}
Der Multiindex $n = (n_1, n_2, n_3, n_4) \in \mathbb{Z}^4$ zählt
die Wicklungen in allen vier Torus-Richtungen.

\subsection{Die $\mathbb{Z}_3$-Orbifold-Projektion}

\textbf{[K]} Die $\mathbb{Z}_3$-Symmetriegruppe wirkt auf $T^4$ durch
den Generator $\tau$, der die Phase $\omega = e^{2\pi i/3}$ einführt.
Der Projektionsoperator auf den $\mathbb{Z}_3$-invarianten Unterraum
ist:
\begin{equation}
  P_{\mathbb{Z}_3} = \frac{1}{3}\sum_{k=0}^{2} \tau^k.
  \label{eq:PZ3}
\end{equation}
Der Hilbertraum des Orbifolds ist:
\begin{equation}
  \mathcal{H}(T^4/\mathbb{Z}_3) = \mathrm{Fix}(P_{\mathbb{Z}_3})
  \subset L^2(T^4).
  \label{eq:Horb}
\end{equation}
Ein Zustand $\psi$ liegt in $\mathcal{H}(T^4/\mathbb{Z}_3)$ genau dann,
wenn $P_{\mathbb{Z}_3}\psi = \psi$, d.h.\ $\tau\psi = \psi$.
\textit{(Korrigiert am 1.~Okt.~2026: vorher $P_{\mathbb{Z}_3} = \tfrac13\sum_k \omega^k\tau^k$ und \glqq $\tau\psi = \omega^{-k}\psi$ für ein $k$\grqq{} -- für $\tau\psi = \lambda\psi$ ist $\tfrac13\sum_k(\omega\lambda)^k\psi$ gleich $\psi$ nur für $\lambda = \omega^{-1}$, sonst $0$; der alte Operator projiziert auf den geladenen Sektor $\tau\psi = \omega^{-1}\psi$, der invariante Projektor ist $\tfrac13\sum_k\tau^k$. Gl.~\eqref{eq:nufix} ist entsprechend mit diesem Projektor zu lesen.)}

\subsection{Fixpunkte des Orbifolds und Neutrino-Lokalisierung}

\textbf{[K]} Die Fixpunkte $x_0$ des Orbifolds erfüllen
$\tau(x_0) = x_0$. Auf $T^4/\mathbb{Z}_3$ gibt es mehrere solcher
Fixpunkte; die für den Neutrinosektor relevanten (Dok.~320) sind:
\begin{equation}
  F_2 = \tfrac{2\pi R}{3}(1,1,1,0),\quad
  F_5 = \tfrac{4\pi R}{3}(1,1,1,0),\quad
  F_7 = \tfrac{4\pi R}{3}(1,1,1,1).
  \label{eq:fixpunkte}
\end{equation}
An einem Fixpunkt $x_0$ gilt die Randbedingung
\begin{equation}
  \psi_\chi(x_0 + y) = \chi \cdot \psi_\chi(x_0 + g_*(y)),
  \label{eq:rbd}
\end{equation}
wobei $\chi$ der Charakter der Darstellung und $g_*$ die
linearisierte $\mathbb{Z}_3$-Wirkung in der Umgebung von $x_0$ ist.

\textbf{[K]} Die drei Neutrino-Flavours sind an diesen Fixpunkten
lokalisiert:
\begin{equation}
  \nu_1 \leftrightarrow F_2,\quad
  \nu_2 \leftrightarrow F_5,\quad
  \nu_3 \leftrightarrow F_7.
  \label{eq:nulok}
\end{equation}
Die unterschiedlichen Fixpunkt-Positionen erzeugen unterschiedliche
Eigenwerte des Massenoperators (Dok.~320, Gl.~6--8).

% =========================================================
\section{Der vollständige Spektraloperator}
% =========================================================

\subsection{Tensor-Produkt-Form von $\hat{F}$}

\textbf{[K]} Der Spektraloperator auf $\mathcal{H}_{\mathrm{FFGFT}}$
hat die Tensor-Produkt-Zerlegung:
\begin{equation}
  \hat{F} = \hat{\Delta}_{T^4} \otimes \hat{S} \otimes \hat{W}
  \label{eq:Ftensor}
\end{equation}
mit:
\begin{itemize}
  \item $\hat{\Delta}_{T^4}$: Laplace-Operator auf $T^4/\mathbb{Z}_3$,
    durch die fraktale Skala $t_0$ diskretisiert.
  \item $\hat{S}$: Spin-Operator auf $\mathcal{H}_{\mathrm{spin}}$
    mit Eigenwerten $s = \pm 1/2$.
  \item $\hat{W}$: Wicklungszahl-Operator auf
    $\mathcal{H}_{\mathrm{flavor}}$ mit Eigenwerten $(n_\theta,n_\phi)$.
\end{itemize}

\textbf{[K]} Der Massenoperator für Fermionen ist der spezialisierte
Sub-Operator:
\begin{equation}
  \hat{M}_f = v \cdot \xi^{\hat{p}} \otimes \hat{r},
  \label{eq:Mf}
\end{equation}
mit dem Exponenten-Operator $\hat{p}$ (Spektrum
$\{3/2, 1, 2/3, 9/4, 2, 9/5\}$) und dem
Wicklungsverhältnis-Operator $\hat{r} = \hat{n}_\phi^2/\hat{n}_\theta$.

\subsection{Spektrale Kommutatoren und Symmetrien}

\textbf{[K]} Die fundamentale Kommutatorbedingung
\begin{equation}
  [\hat{F}, \hat{S}] = 0 \quad \text{auf } \mathcal{H}_{\mathrm{FFGFT}}
  \label{eq:komm}
\end{equation}
drückt aus, dass Spektraloperator und Spin-Operator simultane
Eigenzustände haben -- eine notwendige Bedingung für
wohlbestimmte physikalische Zustände.

Die MASA (maximale abelsche Algebra) der Wicklungszahl-Operatoren:
\begin{equation}
  [\hat{n}_\theta, \hat{n}_\phi] = 0
  \label{eq:MASA}
\end{equation}
erlaubt die gleichzeitige Diagonalisierung und liefert die vollständige
spektrale Kennzeichnung $(n_\theta, n_\phi, p)$ der Teilchenmoden.

% =========================================================
\section{Verallgemeinerte Fourier-Transformation auf $T^4/\mathbb{Z}_3$}
% =========================================================

\subsection{Das GFT-Paradigma: drei Axiome}

\textbf{[B]} Die verallgemeinerte Fourier-Transformation (GFT) für
das FFGFT-Orbifold ist durch drei Axiome definiert:
\begin{enumerate}
  \item \textbf{Invertierbarkeit:} Die GFT ist ein Isomorphismus
    zwischen $L^2(T^4/\mathbb{Z}_3, d\mu_f)$ und dem Spektralraum.
  \item \textbf{Isometrie:} Die GFT erhält das Skalarprodukt
    (verallgemeinerter Parseval-Plancherel):
    $\|\tilde{\psi}\|_{\mathrm{spek}} = \|\psi\|_{\mathcal{H}}$.
  \item \textbf{Spektrale Diagonalisierung:} Der Kern $K_\lambda$
    diagonalisiert $\hat{F}$: $\hat{F}\,K_\lambda = \lambda K_\lambda$.
\end{enumerate}

\subsection{Projektionskern und GFT-Formel}

\textbf{[K]} Der Projektionskern der GFT ist:
\begin{equation}
  K_\lambda^{(p)}(x) = \sum_{n \in \mathbb{Z}^4}
  \frac{\delta_{\lambda,\, \Lambda_n^{(p)}}}{|n|^{D_f/2}}\,
  e^{in\cdot x/R},
  \label{eq:kern}
\end{equation}
mit den verallgemeinerten Eigenwerten
\begin{equation}
  \Lambda_n^{(p)} = \frac{|n|^2}{R^2}\cdot\xi^p.
  \label{eq:lambda_n}
\end{equation}
Die GFT eines Zustands $|\psi\rangle \in \mathcal{H}_{\mathrm{FFGFT}}$:
\begin{equation}
  \tilde{\psi}(\lambda) = \langle K_\lambda^{(p)}, \psi \rangle_{\mathcal{H}}
  = \int_{T^4/\mathbb{Z}_3} \overline{K_\lambda^{(p)}(x)}\,
  \psi(x)\, d\mu_f(x).
  \label{eq:gft}
\end{equation}

\subsection{MASA-Basis und vollständige spektrale Kennzeichnung}

\textbf{[K]} Die MASA $\mathcal{A} = \mathrm{span}\{\hat{n}_\theta,
\hat{n}_\phi, \hat{p}\}$ wird durch drei kommutierende Operatoren
erzeugt. Ihre gemeinsamen Eigenzustände
\begin{equation}
  |n_\theta, n_\phi, p\rangle
  = \frac{1}{2\pi R}\, e^{i(n_\theta\theta + n_\phi\phi)}
  \otimes |p\rangle
  \label{eq:basis}
\end{equation}
sind orthonormal:
\begin{equation}
  \langle n_\theta', n_\phi', p' | n_\theta, n_\phi, p \rangle
  = \delta_{n_\theta n_\theta'}\,\delta_{n_\phi n_\phi'}\,
  \delta_{pp'}
\end{equation}
und bilden eine vollständige Basis für
$\mathcal{H}_{\mathrm{flavor}}$.

% =========================================================
\section{Spektrale Herleitung der Teilchenmassen}
\label{sec:massen}
% =========================================================

\subsection{Geladene Leptonen als Torus-Anregungen}

\textbf{[K]} Die Massenoperator-Eigenwertgleichung
\begin{equation}
  \hat{M}_{e,\mu,\tau}\,|n_\theta, n_\phi, p\rangle
  = m_i\,|n_\theta, n_\phi, p\rangle
  \label{eq:meigen}
\end{equation}
liefert mit Gl.~\eqref{eq:Mf}:
\begin{equation}
  m_i = \frac{n_\phi^2}{n_\theta} \cdot \xi^{p_i} \cdot v.
  \label{eq:masse_eigen}
\end{equation}
Die drei geladenen Leptonen entsprechen drei diskreten
Eigenzuständen:
\begin{center}
\begin{tabular}{lcccc}
\toprule
Teilchen & $n_\theta$ & $n_\phi$ & $p_i$ & $m_i$ \\
\midrule
Elektron $e$  & 3 & 2 & $3/2$ & $(4/3)\xi^{3/2}v = 0{,}505$~MeV \\
Myon $\mu$    & 5 & 4 & $1$   & $(16/5)\xi^1 v = 104{,}96$~MeV \\
Tau $\tau$    & 9 & 5 & $2/3$ & $(25/9)\xi^{2/3}v = 1783{,}4$~MeV \\
\bottomrule
\end{tabular}
\end{center}
\textit{(Korrigiert am 1.~Okt.~2026: vorher Elektron $0{,}511$~MeV -- in der Zeile stand der Messwert statt des Formelwerts; $(4/3)\,\xi^{3/2}\cdot 246$~GeV $= 0{,}5050$~MeV (mit $v = 246{,}22$~GeV: $0{,}5054$~MeV), $-1{,}2\,\%$ neben dem Messwert $0{,}511$~MeV. Die Myon-Zeile $104{,}96$~MeV ist der Formelwert mit $v = 246$~GeV.)}
Dies sind keine freien Parameter, sondern diskrete Spektralwerte des
Massenoperators.

\subsection{Neutrinos als Fixpunkt-Eigenzustände}

\textbf{[K]} Die Neutrinos sind Fixpunkt-Eigenzustände des
$\mathbb{Z}_3$-Projektors:
\begin{equation}
  P_{\mathbb{Z}_3}\,|\nu_i\rangle = |\nu_i\rangle.
  \label{eq:nufix}
\end{equation}
Der Massenoperator auf den Fixpunkt-Zuständen gibt:
\begin{equation}
  m_{\nu_i} = m_e \cdot \xi^{p_i}.
  \label{eq:nu_masse}
\end{equation}
Die Exponenten $p_i$ folgen aus der lokalen fraktalen Dimension
$\dim_{\mathrm{lok}}(F_j)$ des jeweiligen Fixpunkts:
\begin{center}
\begin{tabular}{lcccc}
\toprule
Neutrino & Fixpunkt & $\dim_{\mathrm{lok}}$ & $p_i$ & $m_{\nu_i}$ \\
\midrule
$\nu_1$ & $F_2 = \tfrac{2\pi R}{3}(1,1,1,0)$
  & $2 + 1/4$ & $9/4$ & $0{,}976$~meV \\
$\nu_2$ & $F_5 = \tfrac{4\pi R}{3}(1,1,1,0)$
  & $2$       & $2$   & $9{,}084$~meV \\
$\nu_3$ & $F_7 = \tfrac{4\pi R}{3}(1,1,1,1)$
  & $2 - 1/5$ & $9/5$ & $54{,}11$~meV \\
\bottomrule
\end{tabular}
\end{center}
\textit{(Korrigiert am 1.~Okt.~2026: vorher $\dim_{\mathrm{lok}} = 2 - 1/4$ bei $\nu_1$ -- $2 - 1/4 = 7/4 \neq p_1 = 9/4$; die Masse $m_e\,\xi^{9/4} = 0{,}976$~meV gehört zu $p = 9/4$ (mit $7/4$ wären es $84{,}5$~meV), bei $\nu_2$ und $\nu_3$ gilt $p = \dim_{\mathrm{lok}}$.)}

\subsection{Spektrale Auswahlregel für die Exponentenleiter}

\textbf{[S]} Die Kommutatorbedingung
\begin{equation}
  \hat{F}\cdot\hat{\xi}^{1/3} = \hat{\xi}^{1/3}\cdot\hat{F}
  \label{eq:scaleinv}
\end{equation}
macht das Spektrum invariant unter Skalierung mit $\xi^{1/3}$.
Exponenten derselben Bahn unterscheiden sich um Vielfache von $1/3$:
\begin{equation}
  \Delta p = \frac{k}{3},\quad k \in \mathbb{Z}\quad\text{(innerhalb einer Bahn)},
  \label{eq:leiter}
\end{equation}
so $\mu \to \tau$ ($\Delta p = 1/3 = 1/D$, Dok.~158) und $\nu_2$ auf
derselben Bahn. Die Exponentenleiter umfasst mehrere Bahnen: $p_e$,
$p_{\nu_1}$ und $p_{\nu_3}$ liegen nicht auf dem $1/3$-Raster der
Myon-Bahn. Gl.~\eqref{eq:leiter} beschreibt damit eine
\emph{Bahnstruktur} des Operators $\hat{F}$; die Kommutatorbedingung
Gl.~\eqref{eq:scaleinv} ist nicht hergeleitet.
\textit{(Korrigiert am 1.~Okt.~2026: vorher [K] \glqq Die Exponenten $p_i$ gehorchen der diskreten Auswahlregel $\Delta p = k/3$ \ldots{} ein Spektraltheorem des Operators $\hat{F}$\grqq{} -- $3/2 - 1 = 1/2$, $9/4 - 2 = 1/4$, $2 - 9/5 = 1/5$ sind keine Vielfachen von $1/3$, nur $1 - 2/3 = 1/3$ passt; die Kommutatorbedingung erzwingt keine Leiter über alle Exponenten.)}

% =========================================================
\section{Eichsymmetrien aus der Spektraltheorie}
% =========================================================

\subsection{$SU(3)_c$ als spektrale Projektion}

\textbf{[K]} Die acht Gluonen sind Nullmoden des Spektraloperators
unter der $\mathbb{Z}_3$-Projektion:
\begin{equation}
  \hat{F}\,|g_a\rangle = 0,\quad a = 1,\ldots,8.
  \label{eq:gluon}
\end{equation}
Die $SU(3)$-Generatoren (Gell-Mann-Matrizen) werden durch
Spektralprojektoren in der Wicklungszahl-Darstellung realisiert:
\begin{equation}
  \hat{T}^a = \sum_n \frac{\lambda_n^a}{|n|^2}\,|n\rangle\langle n|.
  \label{eq:su3}
\end{equation}
\textbf{[B]} Die vollständige algebraische Ableitung der Gell-Mann-Matrizen
aus den Orbifold-Moden ist in Dok.~321 ausgeführt.

\subsection{Elektroschwache Bosonen als spektrale Phasenverschiebung}

\textbf{[B]} Die Massivierung von $W^\pm$ und $Z^0$ erfolgt durch
eine spektrale Verschiebung:
\begin{equation}
  \hat{M}_{W} = \tfrac{1}{2}\,\hat{g} \otimes \hat{v},
  \qquad m_Z = m_W/\cos\theta_W,
  \label{eq:MWZ}
\end{equation}
mit Eigenwerten $m_W = 80{,}38$~GeV und $m_Z = 91{,}19$~GeV (Dok.~320,
Gl.~22); $m_W$ ist dabei Eingabe über das gemessene $g$.
\textit{(Korrigiert am 1.~Okt.~2026: vorher $\hat{M}_{W,Z} = \tfrac12\,\hat{g}\otimes\hat{v}\otimes\hat{\xi}^{1/2}$ -- $\tfrac12\cdot 0{,}652\cdot 246{,}22$~GeV $= 80{,}3$~GeV ohne Zusatzfaktor (so Dok.~320); mit $\xi^{1/2} = 0{,}01155$ ergäben sich $0{,}93$~GeV.)}
Photon und Gluonen bleiben als Null-Moden masselos.

% =========================================================
\section{Fraktale Dimension und $K_{\mathrm{frak}}$}
\label{sec:dim}
% =========================================================

\subsection{Die fraktale Dimension als Spektralgröße}

\textbf{[K]} Die fraktale Dimension $D_f$ des FFGFT-Raumes ist nicht
a priori gegeben, sondern emergiert aus dem Spektrum von $\hat{F}$
über die Minkowski-Dimension des spektralen Maßes:
\begin{equation}
  D_f = \dim_H(\mathrm{supp}(E))
  = \inf\!\left\{ s : \sum_i \lambda_i^s < \infty \right\}.
  \label{eq:Df_def}
\end{equation}

\textbf{[K]} Die exakte FFGFT-Formel:
\begin{equation}
  \boxed{D_f = 3 - \xi = 3 - \frac{4}{30000}
  = \frac{22499}{7500} \approx 2{,}9998\overline{6}}
  \label{eq:Df}
\end{equation}
folgt aus Gl.~\eqref{eq:Df_def} angewendet auf $\hat{F}$ mit dem
fraktalen Maß Gl.~\eqref{eq:muf}.

\textbf{Hinweis:} Dok.~003 nennt zusätzlich den Näherungswert
$D_f^{\mathrm{frak}} = 2{,}94$ (Fraktaldimension der diskreten
Rekursion). Dieser Wert wird ohne Herleitung aus $\xi$ angegeben
und bezeichnet eine andere physikalische Größe -- die effektive
fraktale Dimension im messbaren Bereich, nicht den Grenzwert
$D_f \to 3$ der Sub-Planck-Skala.

\subsection{$K_{\mathrm{frak}}$ als Interface-Faktor}

\textbf{[K]} Dok.~133 leitet den kumulativen RG-Lauf-Faktor her:
\begin{equation}
  K_{\mathrm{frak}} = 1 - 100\xi = 1 - \frac{400}{30000}
  = \frac{74}{75} \approx 0{,}9866\overline{6}.
  \label{eq:Kfrak}
\end{equation}
Dieser Faktor tritt beim Übergang von der geometrischen (bare)
Beschreibung zur gemessenen (SI-)Beschreibung auf:
\begin{equation}
  m_i^{\mathrm{SI}} = K_{\mathrm{frak}} \cdot C_{\mathrm{conv}}
  \cdot m_i^{\mathrm{bare}}.
  \label{eq:Kumrech}
\end{equation}

\textbf{[K]} In Massenverhältnissen kürzt sich $K_{\mathrm{frak}}$
vollständig heraus (Dok.~012):
\begin{equation}
  \frac{m_\mu}{m_e} = \frac{m_\mu^{\mathrm{bare}}}{m_e^{\mathrm{bare}}}.
  \label{eq:verh}
\end{equation}
Für die Massenverhältnisse in Dok.~320 ist $K_{\mathrm{frak}}$
daher irrelevant.

% =========================================================
\section{Vollständige spektrale Korrespondenztabelle}
% =========================================================

\begin{table}[h!]
\centering
\caption{Alle spektralen Zustände der FFGFT und ihre
  Hilbertraum-Komponenten (Dok.~322)}
\label{tab:spektral}
\begin{tabular}{llllll}
\toprule
Sektor & Zustand & Eigenwert & Formel &
  $\mathcal{H}$-Komponente & Status \\
\midrule
Leptonen
  & $|e\rangle$      & $0{,}505$~MeV  & $\frac{4}{3}\xi^{3/2}v$
    & $\mathcal{H}_{\mathrm{fl}}(3,2)$       & [K] \\
  & $|\mu\rangle$    & $104{,}96$~MeV & $\frac{16}{5}\xi^1 v$
    & $\mathcal{H}_{\mathrm{fl}}(5,4)$       & [K] \\
  & $|\tau\rangle$   & $1783{,}4$~MeV & $\frac{25}{9}\xi^{2/3}v$
    & $\mathcal{H}_{\mathrm{fl}}(9,5)$       & [K] \\
\midrule
Neutrinos
  & $|\nu_1\rangle$  & $0{,}976$~meV  & $m_e\xi^{9/4}$
    & $\mathrm{Fix}(P)\cap F_2$              & [K] \\
  & $|\nu_2\rangle$  & $9{,}084$~meV  & $m_e\xi^2$
    & $\mathrm{Fix}(P)\cap F_5$              & [K] \\
  & $|\nu_3\rangle$  & $54{,}11$~meV  & $m_e\xi^{9/5}$
    & $\mathrm{Fix}(P)\cap F_7$              & [K] \\
\midrule
Bosonen
  & $|g_a\rangle$    & $0$            & $\hat{F}|g_a\rangle=0$
    & $\mathrm{Fix}(P_{\mathbb{Z}_3})$       & [K] \\
  & $|W^\pm\rangle$  & $80{,}38$~GeV  & $\frac{1}{2}g_2 v$
    & $\mathcal{H}_{\mathrm{spin}}$          & [B] \\
  & $|Z^0\rangle$    & $91{,}19$~GeV  & $m_W/\!\cos\theta_W$
    & $\mathcal{H}_{\mathrm{spin}}$          & [B] \\
\bottomrule
\end{tabular}
\end{table}
\textit{(Korrigiert am 1.~Okt.~2026: vorher $|e\rangle$-Eigenwert $0{,}511$~MeV -- Formelwert $\tfrac43\xi^{3/2}\cdot 246$~GeV $= 0{,}505$~MeV, Messwert $0{,}511$~MeV, $-1{,}2\,\%$; vgl.\ Leptonentabelle in Abschnitt~\ref{sec:massen}.)}

% =========================================================
\section{Offene Fragen und Ausblick}
% =========================================================

\begin{table}[h]
\centering
\caption{Status der offenen Punkte in Dok.~322}
\label{tab:offen}
\begin{tabular}{lll}
\toprule
Frage & Status & Dok. \\
\midrule
Selbstadjungiertheit von $\hat{F}$
  & [B] auf $P_0$; ganzer Raum [S] offen & Dok.~327, R145 \\
Klassifikation des Defektindex von $\hat{F}$
  & [B] $(0,0)$ auf $P_0$ bzw.\ für $\mathrm{Re}\,\hat{F}$ & Dok.~327 \\
$SU(3)$-Generatoren algebraisch aus Orbifold-Moden
  & [B] ausgeführt & Dok.~321 \\
Weinberg-Winkel $\theta_W$ aus $\xi$
  & [K] Dok.~323 & --- \\
$\Delta m_{32}^2$: Quelle der $+16{,}0\,\%$-Abweichung
  & [S] offen & Dok.~320 \\
RGE-Brücke $\alpha_s(m_\tau) \to \alpha_s(M_Z)$ in FFGFT
  & [S] offen & Dok.~160 \\
Verbindung zur nichtkommutativen Geometrie (Connes)
  & [S] perspektivisch & --- \\
Experimentelle Testbarkeit: KATRIN, PTOLEMY
  & [K] messbar & --- \\
\bottomrule
\end{tabular}
\end{table}

\begin{keyresult}[Gesamtstatus Dok.~322]
Die Hilbertraum-Einbettung der FFGFT ist strukturell vollständig:
$\hat{F}$ ist definiert, der Zustandsraum Gl.~\eqref{eq:Htensor}
ist konstruiert, alle Massenformeln aus Dok.~320 sind als
Eigenwertgleichungen identifiziert, und die Selbstadjungiertheit
von $\hat{F}$ ist auf dem $\mathbb{Z}_3$-invarianten Sektor $P_0$ bewiesen (Dok.~327, R86, R145): $\xi = \lambda_{\min}(\hat{F}_{D_4})$
ist damit als Spektralwert wohldefiniert, sofern der $D_4$-Sektor in $P_0$ liegt.
\textit{(Korrigiert am 1.~Okt.~2026: vorher \glqq die Selbstadjungiertheit von $\hat{F}$ ist bewiesen\grqq{}, in Tab.~\ref{tab:offen} \glqq[B] ausgeführt\grqq{} -- nach Dok.~327, Satz~3 (korrigiert) ist $\hat{F}$ unter der $\mathbb{Z}_3$-Paarung normal; selbstadjungiert ist es auf $P_0$, auf dem ganzen Raum nur $\mathrm{Re}\,\hat{F}$ (Setzung).)}
Das Teilchenspektrum folgt aus der Spektraltheorie, nicht aus
freien Parametern.
\end{keyresult}
